% Part: first-order-logic
% Chapter: axiomatic-deduction
% Section: provability-consistency

\documentclass[../../../include/open-logic-section]{subfiles}

\begin{document}

\iftag{FOL}
      {\olfileid[fr]{fol}{axd}{prv}}
      {\olfileid[fr]{pl}{axd}{prv}}

\olsection{\usetoken{S}{derivability} et cohérence}

Nous allons maintenant établir plusieurs propriétés de la relation de
!!{derivability}. Elles présentent un intérêt propre, et chacune
jouera un rôle dans la preuve du théorème de complétude.

\begin{prop}\ollabel{prop:provability-contr}
Si $\Gamma \Proves !A$ et si $\Gamma \cup \{!A\}$ est incohérent,
alors $\Gamma$ est incohérent.
\end{prop}

\begin{proof}
Si $\Gamma \cup \{!A\}$ est incohérent, alors
$\Gamma \cup \{!A\} \Proves \lfalse$. D'après la
\olref[ptn]{prop:reflexivity}, $\Gamma \Proves !B$ pour tout
$!B \in \Gamma$. Comme $\Gamma \Proves !A$ par hypothèse, nous avons
donc $\Gamma \Proves !B$ pour tout
$!B \in \Gamma \cup \{!A\}$. La \olref[ptn]{prop:transitivity} donne
alors $\Gamma \Proves \lfalse$, c'est-à-dire que $\Gamma$ est
incohérent.
\end{proof}

\begin{prop}
\ollabel{prop:prov-incons}
$\Gamma \Proves !A$ si et seulement si
$\Gamma \cup \{\lnot !A\}$ est incohérent.
\end{prop}

\begin{proof}
Supposons d'abord $\Gamma \Proves !A$. Par la
\olref[ptn]{prop:monotonicity},
$\Gamma \cup \{\lnot !A\} \Proves !A$. Par la
\olref[ptn]{prop:reflexivity},
$\Gamma \cup \{\lnot !A\} \Proves \lnot !A$. Nous avons également
$\Proves \lnot !A \lif (!A \lif \lfalse)$ d'après
l'\olref[prp]{ax:lnot2}. Deux applications de la
\olref[ded]{prop:mp} donnent donc
$\Gamma \cup \{\lnot !A\} \Proves \lfalse$.

Supposons maintenant que $\Gamma \cup \{\lnot !A\}$ soit incohérent,
c'est-à-dire que
$\Gamma \cup \{\lnot !A\} \Proves \lfalse$. D'après le théorème de
la déduction, $\Gamma \Proves \lnot !A \lif \lfalse$. En outre,
$\Gamma \Proves (\lnot !A \lif \lfalse) \lif \lnot\lnot !A$ d'après
l'\olref[prp]{ax:lfalse2}; ainsi,
$\Gamma \Proves \lnot\lnot !A$ par la \olref[ded]{prop:mp}. Enfin,
comme $\Gamma \Proves \lnot\lnot !A \lif !A$ d'après
l'\olref[prp]{ax:dne}, une nouvelle application de la
\olref[ded]{prop:mp} donne $\Gamma \Proves !A$.
\end{proof}

\begin{prob}
Démontrer que $\Gamma \Proves \lnot !A$ si et seulement si
$\Gamma \cup \{!A\}$ est incohérent.
\end{prob}

\begin{prop}\ollabel{prop:explicit-inc}
Si $\Gamma \Proves !A$ et $\lnot !A \in \Gamma$, alors $\Gamma$ est
incohérent.
\end{prop}

\begin{proof}
D'après l'\olref[prp]{ax:lnot2},
$\Gamma \Proves \lnot !A \lif (!A \lif \lfalse)$. Deux applications
de la \olref[ded]{prop:mp} donnent $\Gamma \Proves \lfalse$.
\end{proof}

\begin{prop}\ollabel{prop:provability-exhaustive}
Si $\Gamma \cup \{!A\}$ et $\Gamma \cup \{\lnot !A\}$ sont tous deux
incohérents, alors $\Gamma$ est incohérent.
\end{prop}

\begin{proof}
Exercice.
\end{proof}

\tagprob{FOL}
\begin{prob}
Démontrer la \olref[fol][axd][prv]{prop:provability-exhaustive}.
\end{prob}
\tagendprob

\tagprob{notFOL}
\begin{prob}
Démontrer la \olref[pl][axd][prv]{prop:provability-exhaustive}.
\end{prob}
\tagendprob

\end{document}
